\documentclass[../DoAn.tex]{subfiles}
\begin{document}

Chương này tổng kết những kết quả mà đề tài đã đạt được, chỉ ra các hạn chế còn tồn tại và đề xuất các hướng phát triển trong tương lai. Phần \ref{sec:ketluan} đúc kết các đóng góp và kết quả chính. Phần \ref{sec:huongphattrien} trình bày những hướng hoàn thiện và mở rộng tiếp theo của hệ thống.

\section{Kết luận}
\label{sec:ketluan}

Đề tài đã xây dựng hoàn chỉnh một hệ thống xác thực hành vi liên tục trên nền tảng Android, đi từ khâu thu thập dữ liệu, xây dựng mô hình, đánh giá thực nghiệm đến triển khai trên thiết bị. So với các mục tiêu đặt ra ở Chương 1, đề tài đã hoàn thành được những nội dung chính sau đây.

Thứ nhất, đề tài đã xây dựng được một ứng dụng thu thập dữ liệu chuyên dụng và qua đó tạo lập một bộ dữ liệu hành vi thực tế gồm 26 người tham gia, thu trên nhiều dòng thiết bị khác nhau trong điều kiện sử dụng tự nhiên. Việc thu dữ liệu đa thiết bị theo giao thức chuẩn hóa giúp bộ dữ liệu phản ánh sát điều kiện triển khai thực tế hơn so với dữ liệu thu trên một vài dòng máy đồng nhất.

Thứ hai, đề tài đã xây dựng và so sánh có hệ thống ba kiến trúc backbone gồm CNN 1D, CNN-LSTM và CNN-BiLSTM trên hai kịch bản nhận dạng tập đóng và tập mở, đồng thời khảo sát đóng góp của nhánh tương tác màn hình qua cơ chế dung hợp. Ngoài việc so sánh kiến trúc, đề tài còn khảo sát một cách có hệ thống các hàm quyết định chuyển embedding thành điểm tin cậy trên cùng bộ mã hóa, và phát hiện rằng dưới cùng giao thức kiểm thử chéo đáng tin, hàm chuẩn hóa điểm theo nhóm nền cải thiện hiệu năng tập mở so với hàm trung bình cosine cơ sở. Phân tích này cũng cho thấy phần lớn mức lỗi cao quan sát được ở đánh giá ban đầu bắt nguồn từ việc chỉ dựa trên một cụm người lạ cố định chứ không phải từ giới hạn của bản thân bộ mã hóa. Kết quả cho thấy nhánh cảm biến quán tính với kiến trúc CNN 1D cho hiệu năng tốt nhất ở kịch bản tập mở, trong khi việc bổ sung nhánh tương tác qua dung hợp tuyến tính chưa cải thiện hiệu năng trong điều kiện dữ liệu hiện có. Trên cơ sở đó, đề tài lựa chọn cấu hình triển khai gồm CNN 1D, nhánh quán tính, cấu hình toàn bộ ba hoạt động và hàm quyết định chuẩn hóa điểm; cấu hình này đạt tỉ lệ lỗi cân bằng khoảng 0,3\% ở kịch bản tập đóng và khoảng $1,4\% \pm 2,4\%$ ở kịch bản tập mở qua kiểm thử chéo leave-users-out sáu vòng.

Thứ ba, đề tài đã thiết kế một hàm quyết định nhận dạng tập mở dựa trên đối sánh embedding với tập anchor của chủ sở hữu, cho phép đăng ký người dùng mới mà không cần huấn luyện lại mô hình, phù hợp với việc phân phối một mô hình dùng chung trên thiết bị. Toàn bộ đường suy luận được triển khai trực tiếp trên Android thông qua TensorFlow Lite, bảo đảm quyền riêng tư và khả năng hoạt động ngoại tuyến.

Thứ tư, đề tài hiện thực hóa một cơ chế dự phòng đơn giản bằng cử chỉ lắc nhằm hoàn tất vòng xác thực: khi phát hiện người lạ, hệ thống khôi phục phiên ngay trong cùng phương thức chuyển động thay vì quay về phương thức bảo mật truyền thống. Đây là một lựa chọn thiết kế ở mức hệ thống nhằm khép kín chu trình; bản thân cơ chế còn đơn giản và là một hướng có thể phát triển thêm.

Bên cạnh các kết quả đạt được, đề tài còn một số hạn chế cần được nhìn nhận thẳng thắn. Mặc dù việc khảo sát hàm quyết định đã kéo tỉ lệ lỗi cân bằng tập mở xuống mức thấp trong điều kiện thực nghiệm, kết quả này được đo ở mức phiên trên bộ dữ liệu 26 người với số lần thử của người lạ trong mỗi vòng còn nhỏ, nên độ lệch chuẩn vẫn ở mức vài phần trăm và hiệu năng chưa thực sự đồng đều giữa các cá nhân; do đó cần được hiểu là bằng chứng hứa hẹn chứ chưa phải khẳng định về một hệ thống bảo mật đã sẵn sàng. Ngoài ra, nhánh tương tác màn hình chưa phát huy hiệu quả như kỳ vọng và chưa đóng góp vào cấu hình cuối; điều này nhiều khả năng bắt nguồn từ việc lượng dữ liệu touch còn hạn chế và quy trình trích xuất đặc trưng, tiền xử lý cho nhánh này chưa được tinh chỉnh ngang với nhánh quán tính, chứ chưa hẳn do bản chất tín hiệu tương tác, nên cơ chế dung hợp giữa hai nhánh chưa đạt như kỳ vọng. Cuối cùng, cơ chế cập nhật hồ sơ thích ứng tuy đã được hiện thực hóa nhưng chưa được đánh giá định lượng. Quan trọng hơn về mặt bảo mật, hệ thống mới chỉ được đánh giá trước người lạ \textit{zero-effort} (người khác sử dụng thiết bị một cách tự nhiên, không cố bắt chước chủ máy); độ bền trước các tấn công chủ động như bắt chước hành vi hay phát lại tín hiệu cảm biến, cũng như độ an toàn của mã cử chỉ lắc trước quan sát trộm, chưa được khảo sát.

\section{Hướng phát triển trong tương lai}
\label{sec:huongphattrien}

Từ những hạn chế nêu trên, đề tài đề xuất một số hướng phát triển tiếp theo.

Hướng quan trọng nhất là nâng cao khả năng tổng quát hóa của mô hình với người dùng chưa biết nhằm giảm hơn nữa tỉ lệ lỗi và phương sai trong kịch bản tập mở. Việc áp dụng các phương pháp học biểu diễn hướng tới phân biệt danh tính, chẳng hạn học tương phản hoặc các hàm mất mát dựa trên biên góc như ArcFace~\cite{deng2019arcface}, có thể giúp không gian embedding tách biệt các cá nhân tốt hơn và cải thiện khả năng phát hiện người lạ; do quyết định của hệ thống dựa trên độ tương đồng cosine, hướng này phù hợp về mặt hình học và có thể tích hợp mà vẫn giữ nguyên đường triển khai hiện tại.

Hướng thứ hai là cải thiện chất lượng và quy mô dữ liệu. Việc mở rộng số lượng người tham gia, tăng độ đa dạng về thiết bị và bối cảnh sử dụng được kỳ vọng giúp mô hình ổn định hơn giữa các thiết bị khác nhau. Việc tăng quy mô người tham gia còn trực tiếp làm giảm phương sai của ước lượng trong kịch bản tập mở, qua đó cho phép khẳng định mức lỗi một cách chắc chắn hơn.

Hướng thứ ba là hoàn thiện nhánh tương tác màn hình và cơ chế dung hợp. Như đã phân tích ở Chương 5, hiệu năng hạn chế của nhánh tương tác nhiều khả năng đến từ dữ liệu và tiền xử lý hơn là từ bản chất tín hiệu, nên hướng cải thiện trước hết là thu thêm dữ liệu touch và rà soát lại bước trích xuất đặc trưng. Thay cho cách trích đặc trưng thủ công và dung hợp tuyến tính hiện tại, có thể nghiên cứu các phương pháp học đặc trưng trực tiếp từ dữ liệu tương tác và các chiến lược dung hợp thích ứng theo ngữ cảnh, nhằm khai thác hiệu quả hơn tính bổ sung giữa hai nguồn dữ liệu.

Hướng thứ tư là đánh giá toàn diện hơn các thành phần của bản triển khai. Cụ thể, cần đánh giá định lượng tác động của cơ chế cập nhật hồ sơ thích ứng tới độ chính xác theo thời gian, khảo sát ảnh hưởng của kỹ thuật lượng tử hóa mô hình tới chất lượng embedding nhằm giảm thêm dung lượng và độ trễ, đồng thời đo lường hiệu năng thực tế trên thiết bị như độ trễ suy luận và mức tiêu thụ pin trong điều kiện sử dụng dài hạn.

Cuối cùng, để tiến gần hơn tới một sản phẩm hoàn chỉnh, hệ thống cần được kiểm thử trên quy mô người dùng lớn hơn trong thời gian dài, bổ sung đánh giá trước mô hình kẻ tấn công chủ động - gồm tấn công bắt chước hành vi, phát lại dữ liệu cảm biến, và quan sát trộm cử chỉ dự phòng - kèm phân tích định lượng entropy và khả năng bị quan sát của mã cử chỉ lắc, và hoàn thiện trải nghiệm người dùng của cơ chế xác thực dự phòng. Những hướng phát triển này sẽ giúp thu hẹp khoảng cách giữa một hệ thống nghiên cứu và một giải pháp xác thực hành vi liên tục có thể triển khai rộng rãi trong thực tế.

\end{document}